% !TEX program = pdflatex
\documentclass[11pt,a4paper]{article}

\setlength{\parindent}{0pt}

\usepackage[margin=1in]{geometry}
\usepackage[T1]{fontenc}
\usepackage[utf8]{inputenc}
\usepackage{lmodern}
\usepackage{microtype}
\usepackage{booktabs}
\usepackage{array}
\usepackage{amssymb}
\usepackage{enumitem}
\usepackage{xcolor}
\usepackage{hyperref}
\usepackage{titlesec}

\hypersetup{
  colorlinks=true,
  linkcolor=blue!50!black,
  urlcolor=blue!50!black,
  citecolor=blue!50!black
}

\titleformat{\section}{\large\bfseries}{\thesection}{0.6em}{}
\titleformat{\subsection}{\normalsize\bfseries}{\thesubsection}{0.6em}{}

\newcolumntype{P}[1]{>{\raggedright\arraybackslash}p{#1}}

\title{\textbf{Revision Strategy after FADO's First Submission Reviews}}
\author{}
\date{September 2026}

\begin{document}

\maketitle

\section{Overall Assessment}

Strengths:
\begin{itemize}[leftmargin=*]
  \item the problem is practically significant;
  \item the manuscript is generally well organized and clearly written;
  \item the method consistently improves fairness metrics while
    maintaining competitive predictive performance;
  \item the use of 30 random seeds and paired Wilcoxon tests with
    Holm correction is rigorous;
  \item the implementation and experimental configuration are reproducible.
\end{itemize}

The main disagreement is therefore concentrated around
\textbf{novelty, technical justification, and experimental breadth},
rather than a fundamental problem with the research question.

\section{Priority Revision Plan}

\begin{table}[htbp]
  \centering
  \small
  \begin{tabular}{P{0.10\textwidth} P{0.27\textwidth} P{0.50\textwidth}}
    \toprule
    \textbf{Priority} & \textbf{Issue} & \textbf{Why it matters} \\
    \midrule
    1 & Make the novelty of FADO unmistakable & Reviewers 2 and 3
    argue that FADO may be an incremental extension of ARANYANI. Kind
    of critical \\
    2 & Add a comprehensive ablation study & Reviewers explicitly ask
    which individual components are responsible for the gains. This
    is expected \\
    3 & Expand drift experiments & The COMPAS drift is considered
    artificial, and reviewers request gradual, recurring, and
    multiple drift scenarios. Also already discussed as a limitation
    and already have an idea of how to address it. \\
    4 & Justify the reaction controller & The learning-rate,
    temperature, thresholds, and multipliers currently appear
    heuristic. This is essentially just an ablation of the
    controler's components.\\
    5 & Address EO optimization directly & The paper derives EO
    regularization but does not use it for training because of
    computational cost. \\
    6 & Reconsider hyperparameter tuning & Offline NSGA-II tuning may
    appear inconsistent with the fully online motivation. We are not even using NSGA-II as the model is too heavy \\
    7 & Add computational overhead analysis & Reviewers want evidence
    that the method remains practical for streaming systems. Required math that will improve the theoretical foundation. \\
    8 & Discuss interpretability & Oblique decision forests create an
    important limitation for high-stakes applications. Also must do something about it, im just not sure how yet.\\
    \bottomrule
  \end{tabular}
  \caption{Recommended revision priorities.}
\end{table}

\subsection{A new set of experiments to establish novelty}

Compare:
\begin{enumerate}[leftmargin=*]
  \item ARANYANI;
  \item ARANYANI + ADWIN;
  \item ARANYANI + ADWIN + learning-rate adaptation;
  \item ARANYANI + ADWIN + routing-temperature adaptation;
  \item ARANYANI + ADWIN + both adaptations;
  \item Full FADO.
\end{enumerate}

This experiment directly answers the strongest novelty criticism from
Reviewers 2 and 3.

\section{Add a Comprehensive Ablation Study}

The absence of ablations is one of the clearest common weaknesses
across the reviews, which was expected by us.

A useful ablation should separately evaluate:
\begin{itemize}[leftmargin=*]
  \item dual-threshold ADWIN;
  \item learning-rate spike;
  \item routing-temperature adjustment;
  \item confidence filtering;
  \item cooldown;
  \item annealing;
  \item fairness monitoring, where applicable.
\end{itemize}

For each variant, report:
\begin{itemize}[leftmargin=*]
  \item predictive accuracy;
  \item Demographic Parity (DP);
  \item Equalized Odds (EO);
  \item fairness recovery time after drift;
  \item number/frequency of controller activations;
  \item computational cost where feasible.
\end{itemize}

The goal is to answer:
\begin{quote}
  \textbf{Which component causes the observed improvement?}
\end{quote}

This converts FADO from a collection of design choices into a
scientifically analyzed method.

\subsection{Routing-Temperature Adaptation}

Reviewer 3 explicitly asks why lowering $\tau$ during drift
contributes to fairness recovery. This should be supported by:
\begin{itemize}[leftmargin=*]
  \item an explanation of how temperature affects routing and adaptation;
  \item an empirical sensitivity analysis;
  \item ideally, a small theoretical argument about the expected
    effect on adaptation.
\end{itemize}

\subsection{Sensitivity Analysis}

Rather than presenting values such as $\beta_w=5$, $\beta_c=10$, or
$\tau_{\mathrm{drift}}=0.1$ as isolated choices, show that the method
is not excessively sensitive to them.

For example, vary:
\[
  \beta \in \{1,2,5,10,20\}
\]
and evaluate several plausible values of $\tau_{\mathrm{drift}}$.

The central question is whether the main conclusions remain stable
across reasonable parameter choices.

\subsection{Stability and Annealing}

Explain why the temporary aggressive adaptation is followed by
annealing and cooldown. If a formal convergence proof is not
feasible, provide empirical evidence that the controller improves
recovery without causing persistent instability.

\section{Add a proper runtime and efficiency analysis}

The paper already has a strong conceptual argument that it avoids
storing and replaying historical samples. Quantify this advantage.

Report, where possible:
\begin{itemize}[leftmargin=*]
  \item training time per instance;
  \item memory consumption;
  \item runtime overhead relative to ARANYANI;
  \item runtime overhead relative to ARANYANI + ADWIN;
  \item frequency and cost of controller activations;
  \item additional cost of EO optimization.
\end{itemize}

This is especially important because practical viability is part of
the paper's streaming-learning motivation.

\section{Discuss Interpretability and High-Stakes Applications}

Reviewer 1 points out that oblique decision forests are less
interpretable than simple axis-aligned decision rules.

The paper does not necessarily need to abandon its high-stakes
motivation. Instead, it should acknowledge that:
\begin{itemize}[leftmargin=*]
  \item oblique splits can reduce interpretability;
  \item fairness does not by itself make a model suitable for
    high-stakes deployment;
  \item transparency, explainability, and legal requirements remain
    separate concerns.
\end{itemize}

This limitation should be presented explicitly in the discussion or
limitations section.

\end{document}
